\begin{Exercise}[label={ex:inverti}]
    Si consideri una lista di interi concatenata.
    Scrivere una funzione che inverte la lista. Due soluzioni sono possibili:
    \begin{itemize}
        \item la funzione alloca in memoria una nuova lista (invertita) e poi (se richiesto) elimina
        la lista originale;
        \item la funzione riusa i nodi della lista per crearne una nuova invertita.
    \end{itemize}
\end{Exercise}

\begin{Answer}[ref={ex:inverti}]
  	\lstinputlisting[style=cstylesmall]{src/inverti.c}
\end{Answer}